\documentclass{exam}
%==================================================================================
% I. Preambulo
%==================================================================================
%----------------------------------------------------------------------------------
% 1.Paquetes
%----------------------------------------------------------------------------------
\usepackage[spanish]{babel}			% escribir en espanol
\usepackage[utf8]{inputenc}	        % escribir en espanol
\usepackage{amsmath}				% ecuaciones y demases
\usepackage{amsthm}					% ecuaciones y demases
\usepackage{amssymb}				% ecuaciones y demases
\usepackage{graphicx}		        % graficos
\usepackage{float}					% manejar la ubicacion de graficos
\usepackage{verbatim}				% Para escribir codigos
\usepackage{url}					% Para escribir direcciones web
\usepackage{subfig}					% Para poner varias figuras en el mismo marco
\usepackage{psfrag}					% Para hacer reemplazos en las figuras
\usepackage{multicol}
\usepackage{multirow}
\usepackage{bigstrut}
\usepackage{color}
\usepackage{tabularx}               % tablas
\usepackage{amstext}
\usepackage[utf8]{inputenc}
\usepackage{graphicx}
\usepackage{adjustbox}
\usepackage{gensymb}
\usepackage{enumerate}
\usepackage{enumitem}        % Enumerar elementos
\usepackage{enumitem}[a)]
\usepackage{enumerate}[a)]
\usepackage{enumerate}[a.]
\usepackage{float}
\usepackage{hyperref}
\usepackage{pgfplots}
\pgfplotsset{compat=1.17}
\usepackage{comment}
\newtheorem{theorem}{Theorem}
%\theoremsymbol{\ensuremath{\blacksquare}}
%\newtheorem*{solution}{Solución:}

\usepackage{xcolor}

% Margenes y formalidades


%----------------------------------------------------------------------------------
% 2.Estilo de la página
%----------------------------------------------------------------------------------

\pagestyle{headandfoot}					% headers y footers
\headrule 								% Linea horizontal bajo el header
\firstpageheader{Universidad Técnica Federico Santa María}{}{Departamento de Ingeniería Comercial}
\runningheader{Universidad Técnica Federico Santa María}{}{Departamento de Ingeniería Comercial}
\footrule
\footer{}{ \thepage}{}
\parindent = 0pt
\renewcommand\partlabel{(\thepartno.)}
\renewcommand\thesubpart{\roman{subpart}}

%----------------------------------------------------------------------------------
% 3.Formato respuesta
%----------------------------------------------------------------------------------

\printanswers

% \noprintanswers

\renewcommand{\solutiontitle}{\noindent\textsf{\textbf{Respuesta}}\par\noindent}
\renewenvironment{TheSolution}{\vspace{\parskip}\leftskip=0pt\rightskip=0pt
\MakeFramed{\advance\hsize-\width}\solutiontitle\sffamily\ignorespaces}{\unskip
\endMakeFramed}

%----------------------------------------------------------------------------------
% 4.Truquillos
%----------------------------------------------------------------------------------

\newcommand{\rea}[1]{\fbox{\parbox{1.00\textwidth}{\smallskip{\textbf{Fundamentos del Reclamo:}} \begin{rm}\\ #1 \end{rm} \bigskip} }}
\newcommand{\reb}[1]{\fbox{\parbox{1.00\textwidth}{\smallskip{\textbf{Respuesta al Reclamo:}} \begin{rm}\\ #1 \end{rm} \bigskip} }}
\newcommand{\nota}[1]{\fbox{\parbox{0.25\textwidth}{\smallskip{\textbf{Nota:}} \begin{rm}\\ #1 \end{rm} \bigskip} }}
\newcommand{\pr}[2]{\frac{\partial #1}{\partial #2}}
\newcommand{\prv}[2]{\partial #1/\partial #2}
\newcommand{\m}[1]{\mathbf{#1}}
\newcommand{\mm}[1]{\mathcal{#1}}
\newcommand{\req}[1]{ecuacin (\ref{#1})}
\newcommand{\rif}[1]{figura (\ref{#1})}
\newcommand{\rec}[1]{seccin (\ref{#1})}
\newcommand{\integral}[1]{\int_{0}^{\infty}{#1 dt}}
\newcommand{\expo}[1]{\mathbf{e}^{#1}}
\newcommand{\eq}[1]{\begin{eqnarray} #1 \end{eqnarray}}

%==================================================================================
% II. Documento
%==================================================================================

\begin{document}
\begin{center}


\LARGE{\textbf{Guía Estudio Nº6}}\\
\LARGE{\textbf{ICS294}}\\
\bigskip
\bigskip


\normalsize{\textsc{Otoño 2025}}

\end{center}

\section{Límite de probabilidad y distribución límite\footnote{Hint: La ley de esperanzas iteradas dice que E(u_ix_i) = E[E(u_ix_i|x_i)] = E[x_i E(u_i|xi)]}} La clave para la consistencia es obtener la probabilidad de dos promedios, espec\'ificamente de $x_i u_i$ y de $x_i^2$, mediante el uso de Leyes de los Grandes N\'umeros (LLN). Para la normalidad asint\'otica por su parte, la clave est\'a en obtener la distribuci\'on en el l\'imite del promedio $x_i u_i$, mediante un Teorema Central del L\'imite (CLT).\\

Distintos supuestos acerca de las propiedades estoc\'asticas de $u_i$ y $x_i$, llevan a distintas propiedades para $x_i u_i$ y $x_i^2$, y por ende al uso de distintos LLN y CLT. En particular tenemos cuando se tiene \textbf{Muestreo Aleatorio Simple} los $(y_i , x_i)$ se extraen aleatoriamente de la poblaci\'on. De esta forma $x_i$ es iid\footnote{Independiente e idénticamente distribuido} , por lo que $x_i^2$ es iid y $x_i u_i$ es iid si $u_i$ es iid.\\

Teniendo esto en cuenta se le solicita obtener el límite de probabilidad y distribución límite de:
    \[\hat{\beta} = \beta + \frac{\frac{1}{N} \sum_{i=1}^{N} x_i u_i}{\frac{1}{N} \sum_{i=1}^{N} x_i^2}\]
cuando los $x_i$ son obtenidos por muestreo aleatorio simple.\\
\begin{comment}
\textbf{Respuesta:}
    Para el límite de probabilidad se tiene que como $x_i$ y $u_i$ son iid, entonces:
    
    \begin{align*}
    \frac{1}{N} \sum_{i=1}^N x_i^2 &\xrightarrow{p} \mathbb{E}[x_i^2] \\
    \frac{1}{N} \sum_{i=1}^N x_i u_i &\xrightarrow{p} \mathbb{E}[x_i u_i] = 0
    \end{align*}
    
    Esto último ocurre por Ley de Esperanzas Iteradas, que indica que 
    
    $$E(u_ix_i) = E[E(u_ix_i|x_i)] = E[x_i E(u_i|xi)]$$
    
    Y como $E(u_i|x_i)=0$, entonces 
    
    $$E(u_ix_i) = 0$$
    
    Por tanto:
    \[\hat{\beta} \xrightarrow{p} \beta + \frac{0}{\mathbb{E}[x_i^2]} = \beta\]
\end{comment}
%\newpage
\section{Estimación y Consistencia}
Existe un modelo descrito por:

$$y_i = x_i \beta + u_i$$

Donde $E(u_i|x_i) = 0$ y $y_i$, $x_i$ y $u_i$ son escalares. Consideremos el siguiente estimador:

$$\Hat{\beta} = \frac{\overline{y}}{\overline{x}} = \frac{\sum_{i=1}^n y_i}{\sum_{i=1}^n x_i}$$

Se asume que $x_i$ y $u_i$ tienen cuatro momentos finitos y que $\{y_i,x_i\}$ son muestras aleatorias \textit{iid}.

\begin{enumerate}
    \item Encuentre $E[\hat{\beta}|X]$
    \item Encuentre $V[\hat{\beta}|X]$
    \item Muestre que $\hat{\beta} \xrightarrow{\enskip p\enskip} \beta$ cuando $n \longrightarrow \infty$
    \item Encuentre la distribución asintótica de $\sqrt{n}(\hat{\beta} - \beta)$ cuando $n \longrightarrow \infty$
\end{enumerate}

\end{document}